\documentclass[10pt,a4paper]{amsart}
\usepackage{amsmath,amssymb,amsthm,graphicx}
% Retained minimal preamble for the Stacks Project chapter more-algebra.tex (More on Algebra).
% The chapter file pinned-original-source.tex carries only \input{preamble} on its first line, so no
% preamble text lives inside the chapter itself.  The Stacks Project uses one single shared preamble for
% every chapter; the snapshot stages it identically as sources/stacks-chow/preamble.tex and
% sources/stacks-etale-cohomology/preamble.tex.  The block below is copied mechanically, line for line,
% from sources/stacks-chow/preamble.tex lines 14-18 (the Xy-pic setup) and 154-217 (the native theorem
% declarations together with the mathematical macros).  Nothing was added, removed or redefined; the
% wrapper class and the amsmath/amssymb/amsthm/graphicx base packages are supplied by the assembler.
% Chapter: more-algebra.tex -- More on Algebra.
% For commutative diagrams we use Xy-pic
\usepackage[all]{xy}

% We use 2cell for 2-commutative diagrams.
\xyoption{2cell}

% Theorem environments.
%
\theoremstyle{plain}
\newtheorem{theorem}[subsection]{Theorem}
\newtheorem{proposition}[subsection]{Proposition}
\newtheorem{lemma}[subsection]{Lemma}

\theoremstyle{definition}
\newtheorem{definition}[subsection]{Definition}
\newtheorem{example}[subsection]{Example}
\newtheorem{exercise}[subsection]{Exercise}
\newtheorem{situation}[subsection]{Situation}

\theoremstyle{remark}
\newtheorem{remark}[subsection]{Remark}
\newtheorem{remarks}[subsection]{Remarks}

\numberwithin{equation}{subsection}

% Macros
%
\def\lim{\mathop{\mathrm{lim}}\nolimits}
\def\colim{\mathop{\mathrm{colim}}\nolimits}
\def\Spec{\mathop{\mathrm{Spec}}}
\def\Hom{\mathop{\mathrm{Hom}}\nolimits}
\def\Ext{\mathop{\mathrm{Ext}}\nolimits}
\def\SheafHom{\mathop{\mathcal{H}\!\mathit{om}}\nolimits}
\def\SheafExt{\mathop{\mathcal{E}\!\mathit{xt}}\nolimits}
\def\Sch{\mathit{Sch}}
\def\Mor{\mathop{\mathrm{Mor}}\nolimits}
\def\Ob{\mathop{\mathrm{Ob}}\nolimits}
\def\Sh{\mathop{\mathit{Sh}}\nolimits}
\def\NL{\mathop{N\!L}\nolimits}
\def\CH{\mathop{\mathrm{CH}}\nolimits}
\def\proetale{{pro\text{-}\acute{e}tale}}
\def\etale{{\acute{e}tale}}
\def\QCoh{\mathit{QCoh}}
\def\Ker{\mathop{\mathrm{Ker}}}
\def\Im{\mathop{\mathrm{Im}}}
\def\Coker{\mathop{\mathrm{Coker}}}
\def\Coim{\mathop{\mathrm{Coim}}}

% Boxtimes
%
\DeclareMathSymbol{\boxtimes}{\mathbin}{AMSa}{"02}

%
% Macros for moduli stacks/spaces
%
\def\QCohstack{\mathcal{QC}\!\mathit{oh}}
\def\Cohstack{\mathcal{C}\!\mathit{oh}}
\def\Spacesstack{\mathcal{S}\!\mathit{paces}}
\def\Quotfunctor{\mathrm{Quot}}
\def\Hilbfunctor{\mathrm{Hilb}}
\def\Curvesstack{\mathcal{C}\!\mathit{urves}}
\def\Polarizedstack{\mathcal{P}\!\mathit{olarized}}
\def\Complexesstack{\mathcal{C}\!\mathit{omplexes}}
% \Pic is the operator that assigns to X its picard group, usage \Pic(X)
% \Picardstack_{X/B} denotes the Picard stack of X over B
% \Picardfunctor_{X/B} denotes the Picard functor of X over B
\def\Pic{\mathop{\mathrm{Pic}}\nolimits}
\def\Picardstack{\mathcal{P}\!\mathit{ic}}
\def\Picardfunctor{\mathrm{Pic}}
\def\Deformationcategory{\mathcal{D}\!\mathit{ef}}

% --- appended 2026-10-12 (Chief): the project own definitions of the reference/slogan/history
% display environments, copied mechanically from the shared Stacks preamble lines 8-12
% (\usepackage{verbatim} plus the three \newenvironment lines, which the chapter preamble keeps
% verbatim from the project preamble). Nothing added, removed or redefined: these environments
% carry no mathematical content and the project itself typesets them as comments. Without these
% three lines a counted statement whose own text uses them fails with Environment slogan undefined.
\usepackage{verbatim}
\newenvironment{reference}{\comment}{\endcomment}
\newenvironment{slogan}{\comment}{\endcomment}
\newenvironment{history}{\comment}{\endcomment}
\usepackage{cleveref}
\crefname{originallemma}{original lemma}{original lemmas}
\Crefname{originallemma}{Original lemma}{Original lemmas}
\makeatletter
\newlabel{lemma-better-cut-complex-in-two}{{0A1U}{source}{}{sourcepointer-1}{}}
\newlabel{lemma-better-cut-complex-in-two@cref}{{[originallemma][0][]0A1U}{[1][0][]source}{}{}{}}
\makeatother
\title{A cohomology group of a pseudo-coherent complex can be isolated as a free direct summand after localizing}
\author{The Stacks Project Authors}
\date{stacks-project:more-algebra.tex:local-snapshot-20261012; GNU FDL 1.2 or later}
\begin{document}
\maketitle
\setcounter{subsection}{1}
\noindent\textit{Editorial note.} All delivered mathematical text is the author's own, copied from the pinned Stacks Project chapter file more-algebra.tex; no statement or proof step was rewritten, reordered or removed. The counted span is the author environment and its own proof environment. References inside the span are handed to the assembler as navigation stubs that carry the canonical tag of the cited tag only; whether each of them is load-bearing is recorded on the stub itself.\par\medskip
\noindent\textit{Allowed standard prerequisites (editorial).}\begin{itemize}
\item A pseudo-coherent complex whose cohomology in the relevant degrees is finite can be cut in two at a degree after localizing at a suitable element, producing a direct sum decomposition of the truncations. invoked at proof lines 19852-19856 and again at proof lines 19857-19862. Supplies both splittings: the first separates the truncation above the degree from the rest, the second separates the cohomology module in that degree.
\end{itemize}\medskip
\ifdefined\hypertarget\hypertarget{sourcepointer-1}{}\fi
\noindent\textit{Source reference.} Original lemma 0A1U: sources/stacks-more-algebra/pinned-original-source.tex, source lines 19744--19744. Invoked twice, at proof lines 19852-19856 for the upper truncation and at proof lines 19857-19861 for the lower one; each application is made to a different object and the conclusions are combined inside the counted span.\par\medskip
\begin{lemma}
\label{lemma-isolate-a-cohomology-group}
Let $R$ be a ring. Let $\mathfrak p \subset R$ be a prime ideal.
Let $K^\bullet$ be a pseudo-coherent complex of $R$-modules.
Assume that for some $i \in \mathbf{Z}$ the maps
$$
H^i(K^\bullet) \otimes_R \kappa(\mathfrak p)
\longrightarrow
H^i(K^\bullet \otimes_R^{\mathbf{L}} \kappa(\mathfrak p))
\quad\text{and}\quad
H^{i - 1}(K^\bullet) \otimes_R \kappa(\mathfrak p)
\longrightarrow
H^{i - 1}(K^\bullet \otimes_R^{\mathbf{L}} \kappa(\mathfrak p))
$$
are surjective. Then there exists an $f \in R$, $f \not \in \mathfrak p$
such that
\begin{enumerate}
\item $\tau_{\geq i + 1}(K^\bullet \otimes_R R_f)$ is a perfect
object of $D(R_f)$ with tor amplitude in $[i + 1, \infty]$,
\item $H^i(K^\bullet)_f$ is a finite free $R_f$-module, and
\item there is a canonical direct sum decomposition
$$
K^\bullet \otimes_R R_f \cong
\tau_{\leq i - 1}(K^\bullet \otimes_R R_f) \oplus
H^i(K^\bullet)_f[-i] \oplus
\tau_{\geq i + 1}(K^\bullet \otimes_R R_f)
$$
in $D(R_f)$.
\end{enumerate}
\end{lemma}

\begin{proof}
We get (1) from Lemma \ref{lemma-better-cut-complex-in-two} as well
as a splitting
$K^\bullet \otimes_R R_f = \tau_{\leq i}K^\bullet \otimes_R R_f \oplus
\tau_{\geq i + 1}K^\bullet \otimes_R R_f$
in $D(R_f)$. Applying Lemma \ref{lemma-better-cut-complex-in-two} once
more to $\tau_{\leq i}K^\bullet \otimes_R R_f$
we obtain (after suitably choosing $f$) a splitting
$\tau_{\leq i}K^\bullet \otimes_R R_f =
\tau_{\leq i - 1}K^\bullet \otimes_R R_f \oplus H^i(K^\bullet)_f$ in $D(R_f)$
as well as the conclusion that $H^i(K)_f$ is a flat perfect module, i.e.,
finite projective.
\end{proof}

\end{document}
